\documentclass[publikasi]{unnes-ta_plus}

% (opsional) Bahasa Indonesia untuk judul "Daftar Isi", dll.
\usepackage[indonesian]{babel}

% Jika ingin menyisipkan naskah artikel dalam bentuk PDF:
\usepackage{pdfpages}

\unnesSetTitle{JUDUL ARTIKEL ANDA}
\unnesSetAuthor{Nama Mahasiswa}
\unnesSetNIM{1234567890}
\unnesSetProdi{Nama Prodi}
\unnesSetFakultas{Nama Fakultas}
\unnesSetDegree{Sarjana ....}
\unnesSetCity{Semarang}
\unnesSetYear{2026}
\unnesSetPembimbing{Nama Pembimbing}{NIP Pembimbing}

\begin{document}

\maketitleunnes
\unnesPersetujuanPembimbing
% Untuk laporan (bukan proposal) biasanya ada pengesahan penguji dan pernyataan
% \unnesPengesahanPenguji{Ketua}{Sekretaris}{Penguji 1}{Penguji 2}
% \unnesPernyataanKeaslian

\tableofcontents
\unnesMainMatter

% ===== Bagian utama: naskah artikel (sesuai template jurnal) =====
% Opsi 1: Anda tulis langsung di sini / \input{artikel.tex}
% Opsi 2: Anda sisipkan PDF naskah artikel:
% \includepdf[pages=-]{naskah_artikel.pdf}

\chapter{Naskah Artikel}
Tuliskan naskah artikel di sini atau sisipkan PDF seperti contoh di atas.

\end{document}
